\chapter{总结}

本作品将两层匿名HIBE与IBS组合为可配置的抗量子身份基匹配加密系统，完成双方身份约束下的发钥、加密、解密、错误输入拒绝和对象恢复。短种子混合封装承载可变长应用数据，一次性公钥承诺连接内部身份签名与完整密文字节认证；独立的原语接口支持两种HIBE与五类IBS的十种配对。

主要优化围绕计算、存储和重复使用展开。Sageloli后端通过多项式表示、逐频率低秩计算和精确代数检查，完成论文PARAM-I、PARAM-II的采样与签名生命周期。SRGG四类对象的无损编码，使三档主状态长度缩小约5.05至5.46倍，签名约2.55至2.65倍。Lai接收会话将固定身份对的委派安排在准备阶段，在已有一套系统的实验中，后续解密中位数降至28.6秒，准备约510秒，额外叶钥约34.7 MiB。

完整系统测量也明确了使用取舍。A在主实验中的重复加密和解密分别约78.5毫秒、152毫秒，适合现场交互；B展示原生模块格IBS与紧凑编码，签名生成是主要时间成本；C展示环结构采样和连续接收优化，首次委派与签名尺寸仍较大。密文组成分析进一步指出，A主要承担固定一次性认证开销，C主要承担大签名与HIBE胶囊开销，为后续通信优化提供了具体方向。

系统已有命令行、真实浏览器操作和另一台主机的独立安装记录。后续工作可围绕C的最大消息与多套系统重复、以及完整密文中占比较大的组成继续改进，使参数、资源和使用场景的选择更加充分。
